\ifdefined\pdfoutput
  \pdfoutput=1
\fi
\documentclass[11pt,reqno]{amsart}

\usepackage[margin=1in]{geometry}
\usepackage[T1]{fontenc}
\usepackage[utf8]{inputenc}
\usepackage{lmodern}
\usepackage{microtype}
\usepackage{amsmath,amssymb,amsthm,mathtools}
\usepackage{booktabs,longtable,array,tabularx}
\usepackage{enumitem}
\usepackage{float}
\usepackage{aliascnt}
\usepackage{xcolor}
\usepackage{tikz}
\usepackage{xurl}
\usepackage{hyperref}
\usepackage[nameinlink,capitalize,noabbrev]{cleveref}
\usetikzlibrary{arrows.meta,positioning,shapes.geometric,calc}

\definecolor{branchblue}{HTML}{315A8A}
\definecolor{branchgreen}{HTML}{2F6F4E}
\definecolor{branchamber}{HTML}{8A5A12}
\definecolor{branchgray}{HTML}{F4F5F7}
\hypersetup{
  colorlinks=true,
  linkcolor=branchblue,
  citecolor=branchblue,
  urlcolor=branchblue,
  pdftitle={Auditable Branch Closure for Human--AI PDE Proofs},
  pdfauthor={Guillem Duran Ballester},
  pdfsubject={A branch-complete methodology for AI-assisted PDE proof development},
  pdfkeywords={artificial intelligence, mathematical reasoning, partial differential equations, proof repair, branch closure, Navier--Stokes equations}
}

\newcolumntype{L}[1]{>{\raggedright\arraybackslash}p{#1}}
\newcolumntype{Y}{>{\raggedright\arraybackslash}X}
\renewcommand{\arraystretch}{1.13}
\setlist{itemsep=2pt,topsep=4pt}

\theoremstyle{plain}
\newtheorem{theorem}{Theorem}[section]
\newaliascnt{proposition}{theorem}
\newtheorem{proposition}[proposition]{Proposition}
\aliascntresetthe{proposition}
\newaliascnt{lemma}{theorem}
\newtheorem{lemma}[lemma]{Lemma}
\aliascntresetthe{lemma}
\newaliascnt{corollary}{theorem}
\newtheorem{corollary}[corollary]{Corollary}
\aliascntresetthe{corollary}
\theoremstyle{definition}
\newaliascnt{definition}{theorem}
\newtheorem{definition}[definition]{Definition}
\aliascntresetthe{definition}
\newaliascnt{protocol}{theorem}
\newtheorem{protocol}[protocol]{Protocol}
\aliascntresetthe{protocol}
\theoremstyle{remark}
\newaliascnt{remark}{theorem}
\newtheorem{remark}[remark]{Remark}
\aliascntresetthe{remark}
\numberwithin{equation}{section}
\crefname{theorem}{theorem}{theorems}
\Crefname{theorem}{Theorem}{Theorems}
\crefname{proposition}{proposition}{propositions}
\Crefname{proposition}{Proposition}{Propositions}
\crefname{lemma}{lemma}{lemmas}
\Crefname{lemma}{Lemma}{Lemmas}
\crefname{corollary}{corollary}{corollaries}
\Crefname{corollary}{Corollary}{Corollaries}
\crefname{definition}{definition}{definitions}
\Crefname{definition}{Definition}{Definitions}
\crefname{remark}{remark}{remarks}
\Crefname{remark}{Remark}{Remarks}
\crefname{protocol}{Protocol}{Protocols}
\Crefname{protocol}{Protocol}{Protocols}

\newcommand{\R}{\mathbb R}
\newcommand{\calH}{\mathcal H}
\newcommand{\calL}{\mathcal L}
\newcommand{\calU}{\mathcal U}
\newcommand{\proj}{\mathbf P}
\newcommand{\curl}{\operatorname{curl}}
\newcommand{\divergence}{\operatorname{div}}
\newcommand{\loc}{\mathrm{loc}}
\newcommand{\sem}[1]{[\![#1]\!]}
\newcommand{\sourcebase}{https://github.com/FragileTech/hypostructure/blob/main/to_formalize}

\tikzset{
  >=Latex,
  box/.style={draw=branchblue,very thick,rounded corners=2pt,fill=branchgray,
    align=center,text width=7.1cm,minimum height=9mm,font=\small},
  dec/.style={diamond,draw=branchamber,very thick,fill=yellow!9,
    aspect=2.25,align=center,text width=3.8cm,inner sep=1.5pt,font=\small},
  term/.style={ellipse,draw=branchgreen,very thick,fill=green!7,
    align=center,text width=4.7cm,minimum height=10mm,font=\small},
  openbranch/.style={draw=red!65!black,thick,dashed,rounded corners=2pt,
    fill=red!3,align=center,text width=3.0cm,minimum height=9mm,font=\small},
  route/.style={draw=none,fill=none,inner sep=0pt},
  flow/.style={->,very thick,branchblue},
  routeflow/.style={->,thick,dashed,branchblue}
}

\title[Auditable Branch Closure for Human--AI PDE Proofs]
{Auditable Branch Closure for Human--AI PDE Proofs}
\author{Guillem Duran Ballester}
\address{Fragile Technologies}
\email{guillem@fragile.tech}
\date{}
\subjclass[2020]{Primary 68T01; Secondary 35Q30, 35B65, 68V15}
\keywords{AI-assisted mathematics, proof repair, branch closure, partial differential equations, Stokes equations, Navier--Stokes equations}

\begin{document}
\raggedbottom

\begin{abstract}
We introduce branch closure, a human--AI methodology for developing long PDE
proofs as finite, auditable case decompositions.  A branch records its local
space--time window, analytic hypotheses, gauge and compactness choices, active
case, goal, and supporting evidence.  Every transition must be both locally
justified and exhaustive.  This yields a simple soundness theorem and, before
all branches close, a precise partial result whose only possible counterexamples
lie in the open branches.

The central repair operation differs from dependency-based invalidation.  If a
step asserting a predicate \(P\) fails, the last verified branch is partitioned
into \(P\) and \(\neg P\).  All subsequent arguments are retained verbatim on
the \(P\)-branch; only the complementary branch is new.  For a conjunctive
claim, ordered first-failure branches identify the first missing hypothesis.
This makes proof repair monotone and gives an AI system a bounded task without
weakening mathematical completeness.

We give a complete finite-energy Stokes example in which a compact-source
pressure split leaves a spatially harmonic tail.  Solenoidal cutoff estimates
regularize the velocity, a local Caccioppoli contradiction excludes the
parasitic harmonic branch, and the ecodexquation closes the pressure tail without a
whole-space pressure or heat-kernel formula.  The protocol is then instantiated
in three companion developments for the three-dimensional Navier--Stokes equations: a local
concentration and Seregin-profile reduction for the Type I regime, a Liouville
closure of the residual Type I profile class, and a local exclusion of the
complementary Type II regime.  Together, these developments exhaust the local
Type I/Type II dichotomy and rule out finite-time singular points for
finite-energy suitable weak solutions on \(\R^3\).
\end{abstract}

\maketitle

\section{Introduction}
\label{sec:introduction}

Long proofs in nonlinear PDE are difficult to develop with language models for
a reason that is only partly linguistic.  Their conclusions depend on choices
that remain implicit in ordinary prose: the cylinder on which an estimate is
valid, the scale and frame, the pressure gauge, the quotient representative,
the topology used at a compactness limit, and the alternative left after an
estimate fails.  A locally plausible paragraph can silently change one of
these choices and thereby change the theorem.

Current AI systems for mathematical reasoning emphasize formal proof search,
autoformalization, retrieval, and repair of rejected proof scripts
\cite{WuEtAl2022,JiangEtAl2022,FirstEtAl2023,LuEtAl2024}.
Recent PDE-facing systems instead translate control requests or generate
numerical solvers \cite{SorocoEtAl2025,LiEtAl2025}.  These are important but
different problems.  The present paper addresses the intermediate activity in
which a mathematician and an AI system jointly construct a rigorous analytic
argument whose decisive estimates may remain in conventional mathematical
language.

Our proposal is to represent the proof as an exhaustive directed acyclic graph
of analytic cases.  Nodes are not fragments of prose.  Each node exposes the
state required to interpret its claim; each edge records one implication or
one exhaustive split; each terminal records a proof of the target on that
case.  The representation is especially natural for concentration compactness,
profile decompositions, local pressure splittings, scale selection, and
rigidity arguments, where the analysis already advances by alternatives
\cite{Lions1984a,Lions1984b,Gerard1998,BahouriGerard1999,Keraani2001}.

The key point is how the graph changes when an intermediate claim is wrong.
Suppose a node was intended to establish properties \(X\) and \(Y\), and every
later step uses only the declared interface \(X\wedge Y\).  Discovering that
the attempted proof does not establish \(Y\) does \emph{not} erase the later
mathematics.  The preceding case is replaced by the exhaustive alternatives
\[
  \neg X,\qquad X\wedge\neg Y,\qquad X\wedge Y.
\]
The existing continuation is attached unchanged to the third alternative.
The first two alternatives are new analytic obligations.  The root theorem is
complete once they close; until then, they describe exactly what remains.  We
call this operation \emph{branch-preserving repair}.

The paper makes four contributions.

\begin{enumerate}
  \item It defines a PDE-specific branch state and an interface discipline for
  human--AI proof development.
  \item It proves soundness, an exact partial-closure statement, and
  branch-preserving and monotone repair results.
  \item It gives an operational division of labor: AI proposes and maintains
  local branches, while the mathematician controls theorem statements,
  admissible imports, and the acceptance of analytic evidence.
  \item It supplies a complete finite-energy Stokes instantiation that closes
  a harmonic pressure tail by local estimates and contradiction, and points to
  three full-scale Navier--Stokes proof graphs using the same reading-guide,
  source-ledger, monotonicity-ledger, and branch-closure discipline.
\end{enumerate}

The method is deliberately lighter than formalization.  It does not certify an
unproved estimate, and graph consistency is not a substitute for analysis.
Its purpose is to make the analytic obligations explicit, preserve valid work
under repair, and expose the exact residual theorem at every stage.

\section{Branch-complete proof states}
\label{sec:method}

\subsection{Semantic setting}

Fix a universe \(\calU\) of admissible mathematical objects.  An element may
be a weak solution, a rescaled sequence, a limiting profile, or a local
space--time configuration.  A theorem has the form
\begin{equation}\label{eq:target}
  A(s)\Longrightarrow T(s),\qquad s\in\calU,
\end{equation}
where \(A\) is the theorem hypothesis and \(T\) its conclusion.  Direct and
contradiction arguments are both included: in the latter case, the active goal
is the emptiness of \(A\wedge\neg T\).

\begin{definition}[PDE branch state]\label{def:branch-state}
A branch state is a record
\[
  B=(W,H,\Gamma,\tau,O,E).
\]
Here \(W\) fixes the local window, scale, and coordinate frame; \(H\) is the
list of analytic hypotheses already available on that window; \(\Gamma\)
records gauge, quotient, normalization, and convergence choices; \(\tau\) is
the active case predicate; \(O\) is the goal on this case; and \(E\) contains
the mathematical evidence accepted for the preceding implications.  Its
semantic region is
\[
  \sem B=\{s\in\calU: W(s)\wedge H(s)\wedge\Gamma(s)\wedge\tau(s)\}.
\]
The fields are mathematical data.  They may be written in prose, equations,
or a table; no particular serialization is part of the method.
\end{definition}

The separation between \(H\), \(\Gamma\), and \(\tau\) matters in PDE.  A
bound can survive rescaling while a gauge choice does not; a sequence can
converge weakly in the energy space without giving the strong convergence
needed for a nonlinear term; and a pressure statement can be meaningful only
modulo functions of time.  Recording these choices prevents a downstream
argument from consuming a stronger statement than its parent produced.

\begin{definition}[Valid transition]\label{def:transition}
A transition from a branch \(B\) to children \(B_1,\ldots,B_m\), written
\[
  B\Longrightarrow B_1\vee\cdots\vee B_m,
\]
is valid if its evidence establishes
\begin{equation}\label{eq:transition-cover}
  \sem{B_i}\subseteq\sem B\quad(1\le i\le m),
  \qquad
  \sem B\subseteq\bigcup_{i=1}^m\sem{B_i}.
\end{equation}
The first condition is refinement and the second is exhaustiveness.  Children
may overlap, although disjoint alternatives are preferable for diagnosis.  A
one-child transition is an implication; a multi-child transition is a case
split.
\end{definition}

Every transition also declares an \emph{output interface}: the facts that its
children may use.  A proof below the transition is interface-respecting when
it uses the transition only through those facts, not through an unstated detail
of the argument that produced them.  This familiar modularity condition is
what makes branch-preserving repair possible.

\begin{definition}[Closed proof graph]\label{def:closed-graph}
A finite rooted directed acyclic graph is branch-complete for
\eqref{eq:target} if:
\begin{enumerate}
  \item \(\{s:A(s)\}\subseteq\sem{B_0}\) for its root \(B_0\);
  \item every nonterminal node has a valid transition to all of its children;
  \item every terminal \(L\) has accepted evidence for
  \(\sem L\subseteq\{s:T(s)\}\).
\end{enumerate}
A terminal may close by a direct estimate, a contradiction, a reduction to an
earlier theorem, or identification with an explicitly classified exceptional
case.
\end{definition}

The graph is a proof presentation, not a second mathematical theory.  Its
benefit is that coverage, local hypotheses, and sources become reviewable
separately from the prose surrounding them.

\subsection{PDE transitions}

The states used in this paper admit five recurring transitions.

\begin{enumerate}
  \item A \emph{localization} replaces a global object by nested cylinders or
  frequency windows and records every cutoff error.
  \item A \emph{decomposition} splits velocity, pressure, profiles, scales, or
  defects into named components and retains the remainder.
  \item An \emph{estimate} closes the component covered by a theorem whose
  hypotheses appear in the current branch.
  \item A \emph{compactness or rigidity step} changes the state only in the
  topology actually obtained and records the property transported to the
  limit.
  \item A \emph{decision} introduces exhaustive positive and negative cases
  for a property not yet known.
\end{enumerate}

For each node a reader should be able to answer six questions: what case is
active, what inputs are available, what new fact is proved, which result proves
it, which alternatives leave the node, and where each alternative closes.  In
the Stokes example these questions are answered by
\Cref{tab:stokes-node-audit}; the larger Navier--Stokes instances use the same
format.

\subsection{Human--AI division of labor}

\begin{protocol}[Branch-closure loop]\label{prot:loop}
Starting from \eqref{eq:target}, the collaborators repeat the following steps.

\begin{enumerate}
  \item The mathematician fixes the theorem, admissible solution class, and
  which external results may be cited.
  \item The AI proposes the next implication or exhaustive split, together
  with its input and output interfaces and a source-level argument.
  \item The mathematician accepts, rejects, or weakens that argument.  Simple
  graph checks may verify numbering, acyclicity, coverage declarations, and
  source links, but acceptance of an analytic estimate remains mathematical
  review.
  \item An accepted transition is added.  A rejected assertion \(P\) is
  replaced by the split \(P\vee\neg P\); the existing continuation is retained
  on \(P\), and work proceeds on \(\neg P\).
  \item The process stops only when every reachable branch has a terminal
  proof of the target.
\end{enumerate}
\end{protocol}

This allocation uses the model where it is effective: proposing decompositions,
searching for missing cases, restating local obligations, tracing citations,
and maintaining a large graph.  It leaves theorem ownership and analytic
judgment with the human author.  In particular, a fluent proof sketch is never
itself evidence that a branch is closed.

\section{Correctness and repair}
\label{sec:correctness}

\begin{theorem}[Soundness]\label{thm:soundness}
If a finite graph is branch-complete for \eqref{eq:target}, then
\(A(s)\Rightarrow T(s)\) for every \(s\in\calU\).
\end{theorem}

\begin{proof}
Fix \(s\) with \(A(s)\).  The root condition gives \(s\in\sem{B_0}\).
At every nonterminal node containing \(s\), exhaustiveness in
\eqref{eq:transition-cover} gives a child containing \(s\).  Finiteness and
acyclicity produce a terminal \(L\) with \(s\in\sem L\).  Terminal closure
then gives \(T(s)\).
\end{proof}

The same argument describes an incomplete proof exactly.

\begin{theorem}[Partial closure]\label{thm:partial}
Suppose every transition in a finite rooted graph is valid, but a set
\(\calL_{\mathrm{open}}\) of reachable leaves has not yet been closed.  If all
other leaves prove \(T\), then
\begin{equation}\label{eq:partial}
  \{s:A(s)\wedge\neg T(s)\}
  \subseteq
  \bigcup_{L\in\calL_{\mathrm{open}}}\sem L.
\end{equation}
Thus the open leaves, rather than an informal list of gaps, are the exact
residual classes in which a counterexample may remain.
\end{theorem}

\begin{proof}
Route a putative \(s\) with \(A(s)\wedge\neg T(s)\) from the root as in the
proof of \Cref{thm:soundness}.  It cannot reach a closed leaf, so it reaches an
open one.
\end{proof}

\subsection{Branch-preserving repair}

Dependency-graph language often suggests that a failed lemma invalidates its
forward dependency closure.  That is not the repair rule here.  The downstream
proofs are conditional arguments, and a failed attempt to establish their
input does not make those conditional arguments false.

\begin{theorem}[Branch-preserving repair]\label{thm:repair}
Let \(B\) be the last verified branch before a transition intended to establish
a predicate \(P\).  Assume every descendant of that transition is
interface-respecting: it uses the transition only through \(P\) and the facts
already present in \(B\).  If the evidence for \(P\) fails, replace the
transition by
\begin{equation}\label{eq:binary-repair}
  B\Longrightarrow B^+\vee B^- ,
  \qquad
  \sem{B^+}=\sem B\cap\{P\},
  \quad
  \sem{B^-}=\sem B\cap\{\neg P\}.
\end{equation}
Then:
\begin{enumerate}
  \item the replacement transition is valid independently of the failed
  evidence;
  \item the entire former descendant graph remains valid, without alteration,
  below \(B^+\);
  \item the only new proof obligation is \(B^-\); and
  \item if the former descendant graph closed \(B^+\), the repaired graph
  satisfies \eqref{eq:partial} with \(B^-\) as its additional residual class.
\end{enumerate}
\end{theorem}

\begin{proof}
The two sets in \eqref{eq:binary-repair} refine \(\sem B\) and cover it by
excluded middle.  On \(B^+\), the predicate \(P\) is part of the case
definition.  Every former descendant therefore receives exactly the interface
it previously consumed, so its proof is unchanged.  No claim has been made on
\(B^-\); it is the new open leaf.  The final assertion is
\Cref{thm:partial}.
\end{proof}

\Cref{fig:branch-repair} shows the repaired local graph.  The descendant
subgraph is not regenerated; it is reattached to the case in which its declared
input \(P\) holds.

\begin{figure}[H]
\centering
\resizebox{0.94\textwidth}{!}{%
\begin{tikzpicture}[node distance=12mm]
  \node[box,text width=2.6cm] (repairB) {last verified branch \(B\)};
  \node[dec,text width=1.7cm,right=of repairB] (repairP) {does \(P\) hold?};
  \node[box,text width=3.0cm,right=18mm of repairP] (repairD)
    {existing descendants \(D\), consuming only \(P\)};
  \node[term,text width=2.0cm,right=of repairD] (repairT)
    {former closed terminals};
  \node[openbranch,below=11mm of repairP] (repairN)
    {new residual \(B\wedge\neg P\)};

  \draw[flow] (repairB) -- (repairP);
  \draw[flow] (repairP) -- node[above,font=\small]{yes} (repairD);
  \draw[flow] (repairD) -- (repairT);
  \draw[flow] (repairP) -- node[right,font=\small]{no} (repairN);
\end{tikzpicture}
}
\caption{Branch-preserving repair.  The positive arm retains every
interface-respecting descendant; the negative arm is the only new analytic
obligation.}
\label{fig:branch-repair}
\end{figure}

\begin{corollary}[Ordered first-failure repair]\label{cor:first-failure}
If the failed node asserted
\(P=X_1\wedge\cdots\wedge X_m\), it may be replaced by the disjoint exhaustive
split
\begin{equation}\label{eq:first-failure}
\begin{aligned}
  &\neg X_1,\\
  &X_1\wedge\neg X_2,\\
  &\hspace{1.2em}\vdots\\
  &X_1\wedge\cdots\wedge X_{m-1}\wedge\neg X_m,\\
  &X_1\wedge\cdots\wedge X_m.
\end{aligned}
\end{equation}
The former descendants remain on the final branch.  Each new branch records
the first absent property and can be analyzed independently.
\end{corollary}

For the motivating two-property example, discovering that \(Y\) is not
established creates the case \(X\wedge\neg Y\); it does not alter any theorem
proved under \(X\wedge Y\).  The additional case \(\neg X\) is retained when
\(X\) was also part of the unproved conjunctive node.  If \(X\) was already
verified earlier, that branch is empty and may be closed immediately.

\begin{theorem}[Monotone closure]\label{thm:monotone}
Assume repairs are made by \Cref{thm:repair} and accepted evidence is changed
only when its own statement or proof is found defective.  Closing or refining
one open branch does not reopen any other closed branch.  If every newly
created complementary branch is eventually closed, the resulting graph proves
the original theorem.
\end{theorem}

\begin{proof}
A repair partitions one semantic region and leaves all other regions and their
evidence unchanged.  Refinement preserves this separation.  When all leaves
are closed, \Cref{thm:soundness} applies.
\end{proof}

The role of dependency information is therefore diagnostic rather than
destructive.  It identifies the interface at which to insert the dichotomy and
the continuation to attach to its positive arm.  It does not mark a forward
set of mathematical statements as invalid.

\section{Specialization to PDE proofs}
\label{sec:pde-specialization}

For a PDE argument, the abstract record in \Cref{def:branch-state} is made
concrete as follows.

\begin{table}[H]
\caption{Minimal information carried by a PDE branch.}
\label{tab:pde-state}
\centering
\small
\begin{tabularx}{\textwidth}{@{}L{0.16\textwidth}Y Y@{}}
\toprule
Field & Typical content & Question answered \\
\midrule
\(W\) & Cylinder, cutoff, scale, frame, time orientation & Where is the claim valid? \\
\(H\) & Equation, solution class, bounds, sign or flux assumptions & Which analytic hypotheses are available? \\
\(\Gamma\) & Pressure gauge, quotient, normalization, convergence topology & Which representative and limiting meaning are fixed? \\
\(\tau\) & Profile type, scale relation, active defect, positive/negative case & Which alternative is being analyzed? \\
\(O\) & Regularity, decay, contradiction, classification, or reduction & What must this branch prove? \\
\(E\) & Lemma, calculation, citation, and exact input/output interface & Why may the reader traverse the incoming edge? \\
\bottomrule
\end{tabularx}
\end{table}

Three constraints are particularly important.

\paragraph{Locality.}
An interior estimate consumes data on an explicitly larger window and produces
a conclusion on an explicitly smaller one.  Cutoff errors and pressure
remainders are children of the same node unless a cited estimate controls them.

\paragraph{Representative stability.}
Statements involving pressure, harmonic gradients, potentials, or quotient
spaces specify the representative that survives the transition.  Equality of
vorticity does not, by itself, identify velocity; weak convergence of pressure
does not choose its time-dependent additive constant.

\paragraph{Topological sufficiency.}
Each limit passage names the convergence needed by the next nonlinear
operation.  A missing strong convergence statement is handled as a negative
branch, not silently upgraded from a weak one.

These constraints turn common review questions into branch questions.  If a
profile decomposition omits a scale relation, split on that relation.  If a
pressure estimate controls only the nonharmonic component, retain the harmonic
component as a branch.  If a rigidity lemma requires \(X\wedge Y\), use the
ordered alternatives in \eqref{eq:first-failure}.  In each case the successful
continuation remains available while the complementary analytic class becomes
the next research target.

\section{Complete Stokes instantiation}
\label{sec:stokes}

The example now isolates the nonlocal pressure issue that occurs in the
Navier--Stokes developments while removing the nonlinear term.  On
\(\R^3\times(0,T)\) consider
\begin{equation}\label{stokes:eq:local}
  \partial_tu-\Delta u+\nabla p=-\divergence F,
  \qquad \divergence u=0,
\end{equation}
where \((\divergence F)_i=\partial_jF_{ij}\).  The pressure is defined modulo
a distribution depending only on time.

\begin{definition}[Finite-energy Stokes state]\label{stokes:def:energy-state}
A finite-energy weak solution of \eqref{stokes:eq:local} satisfies
\[
 u\in L^\infty_{\loc}(0,T;L^2(\R^3))
 \cap L^2_{\loc}(0,T;\dot H^1(\R^3)),
 \qquad p\in\mathcal D'(\R^3\times(0,T)),
\]
and the equation holds in distributions.  We assume
\(F\in C^\infty(\R^3\times(0,T))\), with its derivatives locally square
integrable in time and space.  This is the linear analogue of the
finite-energy class used in the Navier--Stokes applications.
\end{definition}

Fix \(B_r\Subset B_\rho\Subset B_R\), choose
\(\chi\in C_c^\infty(B_R)\) with \(\chi=1\) on \(B_\rho\), and let
\(N(x)=(4\pi|x|)^{-1}\), so \(-\Delta N=\delta_0\).  Define the compact-source
pressure and the pressure tail by
\begin{equation}\label{stokes:eq:pressure-split}
 p_{\mathrm{near}}=\partial_i\partial_jN*(\chi F_{ij}),
 \qquad p_{\mathrm{tail}}=p-p_{\mathrm{near}}-c(t),
\end{equation}
where \(c(t)\) fixes the spatial mean on \(B_\rho\).  Then
\(\Delta p_{\mathrm{tail}}=0\) on \(B_\rho\).  The tail contains all influence
not represented by the compactly supported local source; its regularity will
be proved, not assumed.

\begin{remark}[Why finite energy is essential]\label{stokes:rem:counterexample}
For a nonsmooth scalar function \(a(t)\),
\[
 u(x,t)=a(t)e_1,\qquad p(x,t)=-a'(t)x_1,\qquad F=0
\]
is a local distributional Stokes solution with zero vorticity.  It is excluded
from \Cref{stokes:def:energy-state} unless \(a=0\), because a nonzero spatially
constant velocity is not in \(L^2(\R^3)\).  Thus a local curl argument alone
cannot prove the theorem below; the finite-energy branch must be used.
\end{remark}

\begin{theorem}[Finite-energy Stokes regularity by local pressure-tail closure]
\label{stokes:thm:main}
Let \((u,p,F)\) be a finite-energy Stokes state on
\(\R^3\times(0,T)\).  Then
\[
 u\in C^\infty(\R^3\times(0,T)),
 \qquad \nabla p\in C^\infty(\R^3\times(0,T)).
\]
For every compact cylinder \(Q'\Subset\R^3\times(0,T)\) and every derivative
order, the corresponding norms of \(u\) and \(\nabla p\) are bounded by the
finite-energy norm of \(u\) on a larger time interval and finitely many local
norms of derivatives of \(F\) on a buffered spatial cylinder.  For every split
\eqref{stokes:eq:pressure-split},
\(\nabla p_{\mathrm{tail}}\in C^\infty(B_r\times I_r)\).
\end{theorem}

The proof uses no whole-space pressure formula and no global heat or Oseen
kernel.  Compactly supported divergence corrections remove pressure from the
velocity estimates.  The only possible failure is a finite-energy parasitic
harmonic mode; local Caccioppoli estimates on expanding balls force that mode
to vanish.  Once the velocity is smooth, the equation closes the harmonic
pressure tail locally.

\subsection{How to read the proof architecture}

Node [2] records the local near/tail pressure split without assigning any time
regularity to the harmonic tail.  Curl and local heat estimates give the
smooth vorticity in [3].  Solenoidal cutoffs and annular Bogovski\u\i{}
corrections then yield the exhaustive alternative [4]: either the local
velocity bootstrap closes, or its first failure produces a nonzero entire
finite-energy field with zero divergence and curl.  The failure branch reaches
[6], where an expanding-ball Caccioppoli estimate contradicts its
normalization.  On the surviving branch, [8] uses the original equation to
prove smoothness of the pressure tail, and [9] gives full regularity of the
unmodified velocity and pressure gradient.

\subsection{Diagram map}

\Cref{tab:stokes-diagram-map} records the single diagram, its node range, and
the branch closed by the local contradiction.

\begin{table}[H]
\caption{Map of the Stokes proof-flow diagram.}
\label{tab:stokes-diagram-map}
\centering
\small
\begin{tabularx}{\textwidth}{@{}L{0.18\textwidth}L{0.20\textwidth}Y@{}}
\toprule
Diagram & Nodes & Role \\
\midrule
\Cref{fig:stokes-proof-flow} & [1]--[9] & Split the pressure locally, close
the only parasitic velocity branch by finite-energy contradiction, and recover
the harmonic pressure tail from the equation. \\
\bottomrule
\end{tabularx}
\end{table}

\subsection{Proof-dependency diagram}

\Cref{fig:stokes-proof-flow} displays the local analytic spine and both outputs
of the solenoidal-bootstrap decision.

\begin{figure}[H]
\centering
\resizebox{0.92\textwidth}{!}{%
\begin{tikzpicture}[node distance=6.2mm]
  \node[box] (n1) {\textbf{[1]} Finite-energy whole-space Stokes state and a buffered local cylinder};
  \node[box,below=of n1] (n2) {\textbf{[2]} Compact-source pressure \(p_{\mathrm{near}}\) and spatially harmonic tail \(p_{\mathrm{tail}}\)};
  \node[box,below=of n2] (n3) {\textbf{[3]} Curl equation and interior heat smoothing of \(\omega=\curl u\)};
  \node[dec,below=9mm of n3] (n4) {\textbf{[4]} Do the local solenoidal difference-quotient estimates close?};
  \node[box,below left=12mm and 18mm of n4] (n5) {\textbf{[5]} No: first failure yields a nonzero entire finite-energy field \(H\) with \(\divergence H=\curl H=0\)};
  \node[term,below=of n5] (n6) {\textbf{[6]} Expanding-ball Caccioppoli estimate forces \(H=0\): contradiction};
  \node[box,below right=12mm and 18mm of n4] (n7) {\textbf{[7]} Yes: all local velocity estimates close, hence \(u\in C^\infty\)};
  \node[box,below=of n7] (n8) {\textbf{[8]} The Stokes equation gives \(\nabla p_{\mathrm{tail}}\in C^\infty\)};
  \node[term,below=of n8] (n9) {\textbf{[9]} Full local regularity of raw \(u\) and \(\nabla p\)};

  \draw[flow] (n1) -- (n2);
  \draw[flow] (n2) -- (n3);
  \draw[flow] (n3) -- (n4);
  \draw[flow] (n4) -- node[above left,font=\small]{no} (n5);
  \draw[flow] (n5) -- (n6);
  \draw[flow] (n4) -- node[above right,font=\small]{yes} (n7);
  \draw[flow] (n7) -- (n8);
  \draw[flow] (n8) -- (n9);
\end{tikzpicture}%
}
\caption{Local closure of the Stokes pressure tail.  The only negative branch
is converted into a finite-energy harmonic mode and closed by contradiction;
the positive branch recovers the unmodified pressure from the equation.}
\label{fig:stokes-proof-flow}
\end{figure}

\subsection{Node-by-node audit table}

\Cref{tab:stokes-node-audit} identifies the precise lemma supporting every
state, decision, and terminal closure.

\small
\begin{longtable}{@{}L{0.06\textwidth}L{0.20\textwidth}L{0.25\textwidth}L{0.40\textwidth}@{}}
\caption{Source ledger for \Cref{fig:stokes-proof-flow}.}
\label{tab:stokes-node-audit}\\
\toprule
Node & Input & Output & Mathematical source \\
\midrule
\endfirsthead
\toprule
Node & Input & Output & Mathematical source \\
\midrule
\endhead
[1] & \Cref{stokes:def:energy-state} and a target cylinder & A finite-energy
state with fixed local buffers & Hypotheses of \Cref{stokes:thm:main}.
\tabularnewline\relax
[2] & [1] and a spatial cutoff & \(p=p_{\mathrm{near}}+
p_{\mathrm{tail}}+c(t)\), with \(\Delta p_{\mathrm{tail}}=0\) locally &
\Cref{stokes:lem:pressure-split}. \tabularnewline\relax
[3] & [1] & Smooth local vorticity &
\Cref{stokes:lem:vorticity}. \tabularnewline\relax
[4] & [1], [3], nested cylinders & Exhaustive success/failure alternative for
the local solenoidal bootstrap & \Cref{stokes:lem:local-bootstrap}. \tabularnewline\relax
[5] & Failure at [4] & A normalized nonzero entire \(L^2\) field with zero
divergence and curl & Failure alternative in
\Cref{stokes:lem:local-bootstrap}. \tabularnewline\relax
[6] & [5] & Contradiction & \Cref{stokes:lem:finite-energy-kernel}.
\tabularnewline\relax
[7] & Success at [4] & \(u\in C^\infty\) on the target cylinder & Iteration
in \Cref{stokes:lem:local-bootstrap}. \tabularnewline\relax
[8] & [2], [7] & Smooth harmonic pressure tail &
\Cref{stokes:lem:pressure-closure}. \tabularnewline\relax
[9] & [2], [7], [8] & Smooth raw velocity and pressure gradient & Proof of
\Cref{stokes:thm:main}. \tabularnewline
\bottomrule
\end{longtable}
\normalsize

\subsection{Monotone-hypothesis ledger}

\Cref{tab:stokes-monotone} records the data retained while cylinders shrink
and the failure branch is exhausted over larger balls.

\begin{table}[H]
\caption{Hypotheses retained along the Stokes route.}
\label{tab:stokes-monotone}
\centering
\small
\begin{tabularx}{\textwidth}{@{}L{0.25\textwidth}Y L{0.22\textwidth}@{}}
\toprule
Datum & How it is used & Persistence \\
\midrule
Nested cylinders & Support cutoffs, annular divergence corrections, and
interior estimates & Windows only shrink. \\
Finite-energy bound for \(u\) & Controls every cutoff estimate and excludes
the parasitic kernel & Retained on every branch and every larger ball. \\
Smooth tensor forcing \(F\) & Supplies the local pressure source and the
difference-quotient right-hand sides & Differentiated only to finite order. \\
Distributional Stokes equation and incompressibility & Remove pressure from
solenoidal tests and recover it after velocity closure & Never replaced by a
projected whole-space equation. \\
Pressure modulo functions of time & Makes the conclusion a statement about
\(\nabla p\) & Fixed in \eqref{stokes:eq:pressure-split}. \\
\bottomrule
\end{tabularx}
\end{table}

\subsection{Branch-closure ledger}

\Cref{tab:stokes-closure} verifies that the success and failure alternatives
both reach a proof of the original theorem.

\begin{table}[H]
\caption{Exhaustive closure of the local Stokes bootstrap.}
\label{tab:stokes-closure}
\centering
\small
\begin{tabularx}{\textwidth}{@{}L{0.25\textwidth}L{0.16\textwidth}Y@{}}
\toprule
Case from [4] & Terminal & Closure \\
\midrule
The local estimates fail & [6] & The normalized first-failure mode is entire,
finite-energy, divergence-free, and curl-free; expanding-ball Caccioppoli
forces it to vanish, contradicting normalization. \\
The local estimates close & [9] & Iteration makes \(u\) smooth; the original
equation then makes \(\nabla p_{\mathrm{tail}}\), and hence \(\nabla p\),
smooth on the smaller cylinder. \\
\bottomrule
\end{tabularx}
\end{table}

\subsection{Analytic proof}
\label{sec:stokes-analytic-proof}

All estimates in this subsection are obtained on nested cylinders.  The
whole-space hypothesis supplies only the finite-energy bound that excludes the
parasitic kernel; no whole-space pressure representation or heat kernel is
used.

\begin{lemma}[Localized pressure splitting]
\label{stokes:lem:pressure-split}
Let \(B_r\Subset B_\rho\Subset B_R\), let
\(\chi\in C_c^\infty(B_R)\) equal one on \(B_\rho\), and define
\(p_{\mathrm{near}}\) and \(p_{\mathrm{tail}}\) by
\eqref{stokes:eq:pressure-split}.  Then
\[
 -\Delta p_{\mathrm{near}}
 =\partial_i\partial_j(\chi F_{ij})
 \quad\hbox{on }\R^3,
 \qquad
 \Delta p_{\mathrm{tail}}=0
 \quad\hbox{on }B_\rho.
\]
For every spatial derivative order, the local norms of
\(p_{\mathrm{near}}\) on \(B_\rho\) are bounded by finitely many local norms
of \(F\) on \(B_R\).
\end{lemma}

\begin{proof}
Taking divergence in \eqref{stokes:eq:local} gives
\begin{equation}\label{stokes:eq:pressure-poisson}
 -\Delta p=\partial_i\partial_jF_{ij}.
\end{equation}
Since \(-\Delta N=\delta_0\), the compact-source potential in
\eqref{stokes:eq:pressure-split} satisfies the first displayed identity.
Because \(\chi=1\) on \(B_\rho\), subtraction from
\eqref{stokes:eq:pressure-poisson} gives the second.  Calder\'on--Zygmund and
interior elliptic estimates give the local bounds
\cite{CalderonZygmund1952,Stein1970}.  The construction uses the Newtonian
parametrix only on the compactly supported source \(\chi F\); it makes no
assertion about a whole-space representative of \(p\).
\end{proof}

\begin{lemma}[Vorticity equation and local smoothing]
\label{stokes:lem:vorticity}
For \(\omega=\curl u\),
\begin{equation}\label{stokes:eq:vorticity}
 (\partial_t-\Delta)\omega=-\curl\divergence F
\end{equation}
in distributions.  Hence \(\omega\) is smooth on every compact cylinder in
\(\R^3\times(0,T)\).
\end{lemma}

\begin{proof}
Apply curl to \eqref{stokes:eq:local}; the pressure term vanishes.  Interior
hypoellipticity for the heat operator on nested cylinders gives smoothness
because \(F\) is smooth.  This is a local parametrix argument: a cutoff is
inserted first, and every commutator is supported away from the target
cylinder \cite[Chapter~IV]{LadyzhenskayaSolonnikovUraltseva1968}.
\end{proof}

\begin{lemma}[Finite-energy harmonic kernel]
\label{stokes:lem:finite-energy-kernel}
If \(H\in L^2(\R^3;\R^3)\) satisfies
\[
 \divergence H=0,
 \qquad \curl H=0
 \quad\hbox{in }\mathcal D'(\R^3),
\]
then \(H=0\).
\end{lemma}

\begin{proof}
The vector identity
\(-\Delta H=\curl\curl H-\nabla\divergence H\) shows that every component of
\(H\) is harmonic.  Let \(\eta_R\) be one on \(B_R\), supported in
\(B_{2R}\), and satisfy \(|\nabla\eta_R|\le C/R\).  Testing the harmonic
equation with \(\eta_R^2H\) gives the local Caccioppoli estimate
\[
 \int_{B_R}|\nabla H|^2
 \le \frac{C}{R^2}\int_{B_{2R}}|H|^2
 \le \frac{C}{R^2}\|H\|_{L^2(\R^3)}^2.
\]
Letting \(R\to\infty\) yields \(\nabla H=0\).  An \(L^2(\R^3)\) constant
field is zero.
\end{proof}

\begin{lemma}[Local solenoidal bootstrap]
\label{stokes:lem:local-bootstrap}
Let \((u,p,F)\) satisfy \Cref{stokes:def:energy-state}.  For nested cylinders
\(Q_r=B_r\times J_r\Subset B_R\times J_R=Q_R\) and every integer \(m\ge0\),
the derivatives of \(u\) of parabolic order at most \(m\) satisfy an estimate
of the form
\begin{equation}\label{stokes:eq:local-bootstrap}
 \sum_{2j+|\alpha|\le m}
 \|\partial_t^j\partial_x^\alpha u\|_{L^2(Q_r)}
 \le C_{m,r,R,J_r,J_R}
 \left(
   \|u\|_{L^2(J_R;L^2(\R^3))}
   +\mathcal N_{m,R,J_R}(F)
 \right),
\end{equation}
where \(\mathcal N_{m,R,J_R}(F)\) contains finitely many local \(L^2\)-norms
of derivatives of \(F\) on \(Q_R\).  Consequently
\(u\in C^\infty(Q_r)\).

\end{lemma}

\begin{proof}
Choose spatial and temporal cutoffs supported in \(Q_R\) and equal to one on
an intermediate cylinder.  Spatial difference quotients of the weak equation
cannot be tested directly after multiplication by a cutoff, because the
cutoff destroys incompressibility.  On the annulus where the cutoff varies,
solve
\[
 \divergence w=\nabla\eta\cdot D_hu,
 \qquad w|_{\partial A}=0,
\]
with the Bogovski\u\i{} operator.  The corrected test field
\(\eta D_hu-w\) is compactly supported and divergence-free.  Testing the
Steklov-averaged equation with this field removes the pressure identically.
The local energy estimate, the boundedness of the annular divergence
correction, and Young's inequality give the first spatial difference-quotient
bound.  Repeating the argument for differentiated equations and using
\eqref{stokes:eq:vorticity} supplies all spatial derivatives.  Temporal
difference quotients are then obtained from the same solenoidal weak equation;
the spatial estimates absorb the Laplacian term.  Iteration over strictly
nested cylinders proves \eqref{stokes:eq:local-bootstrap}.  This is the local
cutoff proof of the standard finite-energy Stokes regularity estimate; see
\cite[Chapters~II--III]{Sohr2001}.  In particular, the constants in
\eqref{stokes:eq:local-bootstrap} come from the cutoff calculation itself; no
compactness normalization of global difference quotients is used.
\end{proof}

\begin{lemma}[Local closure of the harmonic pressure tail]
\label{stokes:lem:pressure-closure}
On the cylinders of \Cref{stokes:lem:pressure-split},
\[
 \nabla p_{\mathrm{tail}}
 =-\divergence F-\partial_tu+\Delta u-\nabla p_{\mathrm{near}}.
\]
Consequently \(\nabla p_{\mathrm{tail}}\) is smooth on every smaller
cylinder, with local estimates inherited from
\Cref{stokes:lem:pressure-split,stokes:lem:local-bootstrap}.
\end{lemma}

\begin{proof}
Rearrange the original Stokes equation and subtract
\(\nabla p_{\mathrm{near}}\).  Every term on the right is smooth on the
smaller cylinder.  This determines the gradient of the harmonic tail without
constructing the pressure outside the cutoff ball.  The additive function
\(c(t)\) is immaterial.
\end{proof}

\begin{proof}[Proof of \Cref{stokes:thm:main}]
Fix an arbitrary compact cylinder and two buffered cylinders around it.
\Cref{stokes:lem:local-bootstrap} gives smoothness of the original velocity
\(u\), not merely of a quotient representative.  Independently, if two
whole-space finite-energy representatives had the same divergence and curl,
their difference would be the field excluded by
\Cref{stokes:lem:finite-energy-kernel}; thus the global representative audit has
no residual harmonic branch.

Apply \Cref{stokes:lem:pressure-split} on the buffered spatial balls.  The
compact-source pressure is smooth.  By
\Cref{stokes:lem:pressure-closure}, the harmonic tail has a smooth gradient on
the target cylinder.  Hence \(\nabla p\) is smooth there.  The target cylinder
was arbitrary, which proves the theorem throughout
\(\R^3\times(0,T)\).
\end{proof}


\section{Full-scale Navier--Stokes instantiations}
\label{sec:navier-stokes-instances}

The Stokes graph is small enough to inspect on one page.  The same method is
used in three companion Navier--Stokes developments whose joint regularity
consequence is stated below.  Each opens with a proof-first reading guide, an
acyclic diagram, a node-by-node source ledger, a monotone-hypothesis ledger,
and an exhaustive branch-closure table.  The counts in
\Cref{tab:full-instances} refer to solid numbered nodes and diagram parts in
those source files.

\begin{table}[H]
\caption{Complete large-scale proof-graph instantiations.}
\label{tab:full-instances}
\centering
\small
\begin{tabularx}{\textwidth}{@{}Y L{0.13\textwidth}L{0.15\textwidth}Y@{}}
\toprule
Development & Nodes & Diagram parts & Source \\
\midrule
Navier--Stokes setup and compactness architecture & 39 & 2 &
\href{\sourcebase/proof_setup.tex}{\texttt{proof\_setup.tex}} \\
Type I residual-closure argument & 159 & 9 &
\href{\sourcebase/type_I_residual_closure.tex}{\texttt{type\_I\_residual\_closure.tex}} \\
Type II regularity argument & 135 & 12 &
\href{\sourcebase/type_II_regularity.tex}{\texttt{type\_II\_regularity.tex}} \\
\bottomrule
\end{tabularx}
\end{table}

The first development proves local concentration, the Type I/Type II
dichotomy, and the Seregin-profile reduction for the Type I branch.  The second
proves that the residual Type I profile class is empty.  The third excludes
physical local Type II singular germs.  Their joint conclusion is the
following assembly theorem.

\begin{theorem}[Joint Navier--Stokes regularity consequence]
\label{thm:joint-ns-regularity}
Let \((u,p)\) be a finite-energy suitable weak solution of the
three-dimensional incompressible Navier--Stokes equations
\[
  \partial_tu-\Delta u+(u\cdot\nabla)u+\nabla p=0,
  \qquad \nabla\cdot u=0,
\]
on \(\R^3\times(0,T)\).  No point \(x_*\in\R^3\) is a singular terminal
point at time \(T\).  Since \(T>0\) is arbitrary, the combined result rules
out finite-time singular points.
\end{theorem}

\begin{proof}
Suppose that \(z_*=(x_*,T)\) is singular.  The local concentration theorem in
the setup development produces positive Caffarelli--Kohn--Nirenberg
concentration and the exhaustive alternative between a local pointwise Type I
bound
\[
  \operatorname*{ess\,sup}_{T-\rho^2<t<T}
  \sqrt{T-t}\,\|u(t)\|_{L^\infty(B_\rho(x_*))}<\infty
\]
for some \(\rho>0\), and its failure, the local Type II alternative.  In the
Type I case, the setup development reduces \(z_*\) to a bounded centered
Seregin profile with invariant local \(L^3\)-nonvanishing.  Its explicit
Liouville classes are closed in the setup development, and its remaining class
is empty by the residual-closure development.  The Type I final assembly
therefore makes \(z_*\) regular, a contradiction.  In the complementary case,
\(z_*\) is a physical local Type II singular germ, which is excluded by the
Type II development.  Both alternatives are impossible.
\end{proof}

The node and diagram counts quantify three complete instantiations of the
representation.  The graphs expose profile and scale alternatives, pressure
gauges, compactness topologies, recurrence routes, and terminal closure
sources.  At this scale, branch-preserving repair retains a large verified
continuation on the positive arm of a failed intermediate assertion while
adding its complementary case as a new branch.

The source repository also contains the broader branch-closure development
from which the present PDE specialization is derived
\cite{DuranBallester2026}.  The present paper isolates the human--AI protocol,
the repair theorem, and the PDE state variables needed to apply it.

\section{Discussion}
\label{sec:discussion}

\subsection{What the method guarantees}

If the mathematical evidence on every edge is correct, the graph guarantees
that no admissible case has been omitted.  If some leaves remain open,
\Cref{thm:partial} identifies the exact residual theorem.  If an asserted
property fails, \Cref{thm:repair} preserves the conditional mathematics below
that property and opens its complement.  These guarantees are semantic and do
not depend on a particular model or proof assistant.

The representation also supports useful automated checks: node identifiers can
be compared with ledger rows; terminal shapes can be checked for outgoing
edges; decision nodes can be required to list all alternatives; source links
can be resolved; and state fields can be compared across transitions.  Such
checks detect structural omissions.  They do not decide whether a cited PDE
estimate is applicable.

\subsection{Limitations}

The method does not ensure that the graph is finite before the mathematics is
understood, nor that every complementary branch is tractable.  A poor state
choice can generate unhelpful cases.  Interface-respecting descendants are
essential: if a later proof relied on an unstated detail of a failed argument,
that later proof must first be rewritten against an explicit interface.  The
approach therefore trades some expository overhead for precise repair and
review boundaries.

The Stokes example is linear and its first-failure contradiction is explicit.
The Navier--Stokes instances are substantially larger, but they do not constitute a
controlled user study.  Measuring whether the protocol improves correctness,
review time, or AI contribution quality is future work.  A useful evaluation
would compare ordinary prose development with branch-complete development on
the same PDE problems and record omitted cases, repair cost, and expert review
time.

\section{Conclusion}

Branch closure gives human and AI collaborators a common unit of work: a local
analytic case with a declared interface and an exhaustive continuation.  Its
most important consequence is a different account of failure.  A broken step
does not destroy its descendants.  It induces a dichotomy at the last verified
state; the old proof survives on the positive arm, and the negative arm becomes
the new mathematical problem.  Soundness follows from exhaustive routing, and
partial progress has an exact semantic meaning.

The Stokes example shows why this matters even in a linear PDE: a local
pressure split leaves a harmonic tail, while vorticity smoothing leaves a
parasitic velocity mode.  The finite-energy branch excludes that mode by an
expanding-ball Caccioppoli contradiction, after which the equation closes the
pressure tail locally.  The accompanying Navier--Stokes
developments show that the same discipline scales to proof graphs with more
than one hundred states.  This provides a concrete basis for AI assistance in
PDE proof construction without conflating structured collaboration with
automatic verification.

\section*{Statement on AI assistance}

Language models were used during manuscript editing and for consistency checks.
The author selected the mathematical claims, verified the analytic arguments,
and is responsible for the final text.  No model output is treated as
mathematical evidence merely by virtue of being generated.

\bibliographystyle{amsalpha}
\bibliography{llm_auditable_proof_architecture_draft}

\end{document}
